% SPDX-License-Identifier: MIT
% ----------------------------------------------------------------------------
% Implementation Differential and Provenance Supplement
% Companion to the didactic manual (docs/manual.pdf) of big.bi-supplement.
% Standalone document, independent pagination (bound after the manual).
% Prepared October 2026 by the automated audit agent; author: F. Berthommier.
% ----------------------------------------------------------------------------
\documentclass[11pt,letterpaper]{article}

% ── Packages ──
\usepackage[T1]{fontenc}
\usepackage[utf8]{inputenc}
\usepackage{lmodern}
\usepackage[margin=2.5cm]{geometry}
\usepackage{microtype}
\usepackage{amsmath}
\usepackage{amssymb}
\usepackage{booktabs}
\usepackage{tabularx}
\usepackage{adjustbox}
\usepackage{tcolorbox}
\usepackage{fancyhdr}
\usepackage[colorlinks=true,linkcolor=blue!60!black,citecolor=blue!50!black,urlcolor=blue!70!black,bookmarks=true,bookmarksnumbered=true]{hyperref}

\tcbuselibrary{breakable}

\hypersetup{
    pdftitle={Implementation Differential and Provenance Supplement — big.bi-supplement vs arXiv:2307.02299},
    pdfauthor={F. Berthommier (prepared by automated audit agent)},
    pdfsubject={Differential between the 2023 article parameters and the 2026 Python engine},
    pdfkeywords={articulatory synthesis, Maeda model, VLAM, synthSYL, provenance, reproduction}
}

% ── Page style ──
\setlength{\headheight}{14.5pt}
\pagestyle{fancy}
\fancyhf{}
\fancyhead[L]{\small\itshape Implementation Differential \& Provenance Supplement --- big.bi-supplement}
\fancyfoot[C]{\thepage}
\renewcommand{\headrulewidth}{0.3pt}
\renewcommand{\footrulewidth}{0pt}

% ── Layout guards (no overfull boxes, no orphan lines) ──
\widowpenalty=10000
\clubpenalty=10000
\tolerance=2000
\emergencystretch=3em

% ── Shorthand ──
\newcommand{\code}[1]{\texttt{#1}}
\newcommand{\Mua}{M.u.}

% ── Boxes ──
\newtcolorbox{qbox}[1]{%
    colback=black!4, colframe=black!50, boxrule=0.4pt, arc=1.5pt,
    left=7pt, right=7pt, top=5pt, bottom=5pt,
    title={#1}, fonttitle=\bfseries\small, coltitle=white, colbacktitle=black!75,
    breakable, before upper={\raggedright}}
\newtcolorbox{pastebox}[1]{%
    colback=black!3, colframe=black!60, boxrule=0.7pt, arc=1.5pt,
    left=7pt, right=7pt, top=5pt, bottom=5pt,
    title={#1}, fonttitle=\bfseries\small, coltitle=white, colbacktitle=black!80,
    breakable, before upper={\raggedright}}

\begin{document}

% ════════════════════════════════════════════════════════════════════════════
% Title block (compact; the document is bound after the manual)
% ════════════════════════════════════════════════════════════════════════════
\begin{center}
{\LARGE\bfseries Implementation Differential and Provenance Supplement}\\[5pt]
{\large From the 2023 article (MATLAB era) to the 2026 Python engine}\\[9pt]
{\small Companion to the didactic manual (\code{docs/manual.pdf}) --- repository
\code{big.bi-supplement},\\ commit \code{4a0d76290aee}, DOI 10.5281/zenodo.23239483}\\[4pt]
{\small October 2026}\\[3pt]
{\small\itshape Documentation supplement --- no engine code is changed.}
\end{center}
\vspace{2pt}
\noindent\rule{\textwidth}{0.6pt}
\vspace{4pt}

\tableofcontents
\vspace{6pt}

\section{Purpose and reading guide}

The didactic manual of this repository explains how to reproduce the
demonstrations of Berthommier (2023) with the Python engine shipped in
\code{big.bi-supplement}. This supplement answers the question that precedes
reproduction: \emph{which} parameter values and assembly rules implemented in
the 2026 engine differ from those published in the 2023 article, \emph{where}
each difference lives in the code, \emph{what} it changes numerically in the
simulations, and \emph{why} the published demonstrations are nevertheless
reproduced. It is the documentation-level fix recommended by the independent
compliance audit \hbox{[AUD]}; it proposes and implies no engine code change.

Every factual claim carries a source tag:
\begin{itemize}\itemsep1pt\parsep0pt
  \item \textbf{[A \S x]} --- the 2023 article, F.\,Berthommier,
        ``Why can big.bi be changed to bi.gbi? A mathematical model of
        syllabification and articulatory synthesis'', arXiv:2307.02299
        (Interspeech 2023), section \S x;
  \item \textbf{[M l.n]} --- the didactic manual, \code{docs/manual.tex}
        (29-page PDF at the same commit), line n;
  \item \textbf{[R file:line]} --- the repository engine or scripts at commit
        \code{4a0d76290aee};
  \item \textbf{[AUD]} --- the 104-test compliance audit report (Oct.\ 2026);
  \item \textbf{[D5]} --- the difference-5 effect report (Oct.\ 2026),
        the follow-up experiment summarized in \S5 below;
  \item \textbf{[EXT \S x]} --- the external comparative audit of the
        consonant selection vectors $C \in \{$b, d, g$\}$ in the two original
        engines, October 13, 2026, reference [6].
\end{itemize}

Section \S7 issues three \emph{corrigenda}: two factually wrong statements of
the manual, superseded by the replacement texts given there, and one family of
wrong labels in the demonstration figures generated by the repository's own
figure generator. No other part of the manual is affected.
The intended audience is researchers who reproduce, extend, or cite the
demonstrations.

Two conventions. First, Maeda parameter indices are 1-based in the article and
the manual and 0-based in the code; article $\{1,2,6\}$ = code $\{0,1,5\}$
(Jaw, Body, LipH), and the seven parameters are Jaw, Body, Dorsum, Tip, LipP,
LipH, Hy \hbox{[R \code{synthSYL/constants.py:19}]}. Second, magnitudes are
expressed in \emph{Maeda units} (\Mua), the natural units of the 7-parameter
articulatory flow $P(t)$.

\section{Provenance and versions}

\textbf{Lineage (author-asserted).} The 2023 article's demonstrations were
produced with the author's \textsc{matlab} implementation. The engine shipped
in this repository is Python: it descends from the author's
\code{timit-to-Maeda} package (``the original synthSYL'') and was reworked
during 2026 to produce every demonstration presented here --- videos, HTML
simulations, article-format figures, and the manual itself. The
\textsc{matlab} source is not part of the audited corpus, and this supplement
makes no claim about its internals; everything published is attributed to the
article. The verifiable axis of this document is therefore:
\emph{values and rules published in the article (2023 era)} $\leftrightarrow$
\emph{values and rules implemented in the engine (2026 era)}
\hbox{[AUD; R]}.

\textbf{What ``verbatim from the original code'' means.} The engine registry
opens with this statement:
\begin{qbox}{\code{synthSYL/constants.py:6--8} (verbatim)}
\small
``All magic numbers from the original monolithic plugin are replaced by
descriptively-named constants here. Every value is taken verbatim from the
original code so that numerical identity is preserved.''
\end{qbox}
\noindent ``The original code'' designates the pre-existing Python engine
(\code{timit-to-Maeda}), \emph{not} the article. Registry inheritance therefore
guarantees numerical identity with the Python lineage --- and that lineage does
not coincide everywhere with the values printed in the article. This
ambiguity is one of the gaps the present supplement resolves. Throughout
\S4, every non-article engine value is labelled either
\emph{registry inheritance} (carried over verbatim from
\code{timit-to-Maeda}, per the quote above) or \emph{author ruling (dated)}.
Two rulings are documented in the code: the word-onset anchoring rule,
author-validated 2026-10-06 \hbox{[R \code{synthSYL/gesture.py:399--406}]},
and \code{VALRECT} $= 1.10$, ruled 2026-10-07
\hbox{[R \code{config.py:43--46}]}.

\textbf{The 2026 rework, as documented.} The README section ``Key
Modifications to the Original synthSYL'' lists four items
\hbox{[R README]}: (1)~\code{delta\_o} and \code{delta\_e} added as pipeline
kwargs; (2)~the syllable-onset anchor at a `.' C.C boundary weighted by
\code{COEFCEN} (the previous vowel weighted by the chaining coefficient);
(3)~the Vo formula corrected to $\delta_o\cdot\rho_V$ (matching the article);
(4)~non-initial word onsets anchored on the previous Ve. Items~1 and~3 conform
to the article's formulas \hbox{[A \S2.2]}; item~2 is the inter-syllable
chaining mechanism of the demonstrations; item~4 is difference~5 of \S4 below
--- the README documents the rule and its motivation, but not its divergence
from a literal reading of the article nor the conditions under which the two
readings coincide (\S5).

\textbf{Version pinning.} Repository \code{big.bi-supplement}, commit
\code{4a0d76290aee}, MIT license, DOI 10.5281/zenodo.23239483 \hbox{[R]}.
The manual (\code{docs/manual.pdf}, 29 pages) and all shipped artifacts were
produced at this commit.

\begin{qbox}{Summary statement [AUD]}
\small
The engine faithfully implements the article's equations, but its parameters
are not exactly those of the 2023 article; the differences are catalogued in
\S4, and none of them is an equation-level change.
\end{qbox}

\section{The model as published (2023 reference)}

The model plans formant trajectories directly in the complex plane. Targets
are polar points $(\rho,\theta)$; the \emph{coordination function} (article
Eq.~1) maps a target to the 7 Maeda parameters,
\begin{equation}
P_i \;=\; \Omega_i \;+\; \rho \cdot \Psi_{1,i} \cdot
\cos\!\left(\Psi_{2,i} - \theta\right),
\label{eq:coordination}
\end{equation}
where $\Omega$, $\Psi_1$, $\Psi_2$ are per-parameter calibration constants
(article Table~1) \hbox{[A \S2.1]}. Motion between two targets follows a
circular arc in the plane (article Eqs.~2--3) with curvature $K$ and the
cosine velocity profile $\rho(t)=\cos(\theta/2)$ with $P_{\text{exp}}=1$
\hbox{[A \S2.1]}. Superposition (article Eq.~4) is selection-based: at every
instant each segment activates a consonant set $S_c$ and a vowel set $S_v$ of
the 7 Maeda parameters \hbox{[A \S2.1--2.2]}.

\textbf{Published settings} \hbox{[A \S2.1]}: $K_{\text{voy}} = 30$ (vowel
arcs), $K = 10$ (consonant arcs), $T = 16$ steps per syllabic period in
didactic mode (1 step $=$ 10~ms, so $T = 160$~ms), $\nu = -1$,
$P_{\text{exp}} = 1$.

\textbf{Anchoring vowels and pauses} \hbox{[A \S2.2]}:
\begin{qbox}{Article \S2.2 --- anchoring vowels (Q1, verbatim [A])}
\small\itshape
``When they are absent, they can be far from the neutral vowel (with $\rho=0$
and $P=\Omega$) and they have the same $\theta$ as the closest vowel but
$\rho$ is multiplied by specific coefficients $\delta_o$ and $\delta_e$
(which will be set later). Thus, we have $V_o=(\delta_o.\rho_V,\theta_V)$ and
$V_e=(\delta_e.\rho_V,\theta_V)$.''
\end{qbox}
\begin{qbox}{Article \S2.2 --- pause (Q2, verbatim [A])}
\small\itshape
``the pause is defined as the lack of concatenation and coarticulation between
two syllables. This is consistently defined as a diphthong transition between
the previous $V_e$ and the next $V_o$ with an arbitrary duration $T_p$.''
\end{qbox}
\noindent The article's demonstrations use $\delta_o=\delta_e=\delta$
throughout, with three published settings: $\delta=0.5$ (salient schwa),
$\delta=0.7$ (audible schwa), and $\delta=1$ (fusion, article Figure~4)
\hbox{[A \S2.2, \S4]}:
\begin{qbox}{Article \S4 --- audible schwa (Q3, verbatim [A])}
\small\itshape
``the graphs of /big/ and /bi/ are concatenated with $\delta_o=\delta_e=0.7$
to insert an audible schwa-like vocoid between the two consecutive
consonants.''
\end{qbox}

\textbf{Published consonant positions and selections} \hbox{[A \S3]}:
/b/ at $(\rho_b{=}1,\ \theta_b{=}\pi/3)$ with $S_c=\{1,2,6\}$; /d/ around
$(\rho_d{=}1.2,\ \theta_d{=}3\pi/2)$ with $S_c=\{1,2,3,4\}$; /g/ velar at
$(1.2, \pi/3)$ and palatal at $(1.1,\ 23\pi/12)$ with $S_c=\{1,2,3,4\}$
(the /g/ is velar for back vowels, palatal for front vowels; the /d/ position
depends slightly on the following vowel). Corner vowels are published at
$\rho = 1$, and the calibration method is stated openly: the model's
parameters ``can be adjusted manually by simple hearing'' \hbox{[A \S1]}, and
``[t]he tuning of /b,d,g/ has been realised by hearing at first''
\hbox{[A \S3]}. For consonant clusters:
\begin{qbox}{Article \S3 --- cluster selection (verbatim [A];
checked against the arXiv rendering for this supplement)}
\small\itshape
``[T]he tuning of the six pair clusters of /b,d,g/ follows the same
principle, but there is a special grouping of articulators when /b/ is
involved in the pair. In this case $S_c=\{1,2,3,6\}$ and otherwise this is
$S_c=\{1,2,3,4\}$.''
\end{qbox}

\textbf{Published demonstration} \hbox{[A \S4]}: \code{big.bi} transforms
into \code{bi.gbi} as $\delta \to 1$; the gain of consonant clustering is one
period $T = 160$~ms. The article's supplement demonstrates
\code{ib\,ib} $\to$ \code{bi\,bi}. Article Figure~3 reports nominal locus
slopes $\approx 0.70$ (/b/), $0.55$ (/d/), $2.0$ (/g/ velar), and $0.75$
(/g/ palatal).

\section{The differential}

Table~\ref{tab:master} is the central exhibit of this supplement: one row per
difference between the article-published values/rules and the engine at
commit \code{4a0d76290aee}. Row~1 carries two entries (1 and 1b): both lie in
the Tip row of the coordination matrix. Paragraphs after the table give each
row's origin, effect, and current documentation status. The two manual
statements quoted in \S7 are superseded by the corrigenda there.

\begin{table}[htbp]
\centering
\caption{Master differential --- article (2023) vs engine (2026), commit
\code{4a0d76290aee}. Indices are 1-based in article numbering unless a code
vector is shown. \Mua{} = Maeda units. Code paths are relative to the repository root:
\code{constants.py}, \code{phonology.py}, \code{gesture.py},
\code{trajectory.py} and \code{polar.py} are under \code{synthSYL/};
\code{config.py} is at the root.}
\label{tab:master}
\scriptsize
\begin{adjustbox}{max width=\textwidth,center}
\begin{tabular}{@{}>{\raggedright\arraybackslash}p{0.75cm}
                >{\raggedright\arraybackslash}p{1.45cm}
                >{\raggedright\arraybackslash}p{2.85cm}
                >{\raggedright\arraybackslash}p{3.25cm}
                >{\raggedright\arraybackslash}p{2.15cm}
                >{\raggedright\arraybackslash}p{3.25cm}
                >{\raggedright\arraybackslash}p{1.95cm}@{}}
\toprule
\textbf{\#} & \textbf{Item} & \textbf{Article (2023)} & \textbf{Engine (2026)}
& \textbf{Code [R]} & \textbf{Quantified effect} & \textbf{Doc.\ status} \\
\midrule
1/1b & CO Tip row ($\Omega$, $\Psi_1$) & $\Omega=0$; $\Psi_1=-3$ (Table~1)
& $\Omega=0.5$; $\Psi_1=-2.75$ (registry inheritance)
& constants.py: 28 & up to 0.70 \Mua{} on Tip for /d/, /g/ (arc
/d/$\to$/a/); zero for /b/, /gb/ [AUD] & omitted --- manual states a
different claim (Corr.~1) \\
2 & $\theta_{\text{dental}}$ & $3\pi/2$ (\S3) & $23\pi/16$
($\Delta=\pi/16=11.25^{\circ}$) & constants.py: 46 & reproduction scripts
restore the article value (\S6) & absent \\
3 & Radii & $\rho_b=1$; corner vowels $\rho=1$ (\S1, \S2.1)
& $\rho_b=1.2$; /i/ $=0.9$; /a/ $=0.8$ (recalibration by ear)
& constants.py: 74, 77, 79 & Body peak $+1.50$ vs $+1.25$ at $\rho_b=1.0$
(manual ablations) [M l.870--872, l.1148--1154] & $\rho_b$ documented; vowel
radii not flagged [M l.276--277] \\
4 & Cluster selection $S_c$ & /b/-clusters $\to \{1,2,3,6\}$; otherwise
$\{1,2,3,4\}$ (\S3) & place-index sum ($b{=}0$, $d{=}1$, $g{=}2$): $\le 2
\to \{1,2,3,6\}$; $>2 \to \{1,2,3\}$ (article numbering; verified by
execution)
& phonology.py: 337--341; constants.py: 89--92 & 4 of 9 /b,d,g/ pairs
diverge (dd, dg, gd, gg); /gb/ conforms [AUD] & wrong --- manual documents
the article's rule (Corr.~2); ancestral [EXT] \\
5 & Word-onset anchoring & literal reading: $V_{o,2}=(\delta_o\,\rho_{V_2},
\theta_{V_2})$ (\S2.2) & $V_{o,2}:=(\delta_e\,\rho_{V_1},\theta_{V_1})$
(author ruling 2026-10-06) & gesture.py: 385--432 & 0 (bitwise) on all 6
published configurations; up to 2.0812 \Mua{} outside them [D5] & partial
(README item~4) \\
6 & VALRECT (soft rectification) & not published & 1.10 (historical
timit-to-Maeda value 0.75; author ruling 2026-10-07) & config.py: 43--46 &
locus-slope effect $\le 0.23$ (/g/ velar), $\le 0.014$ (/b/, /d/) [AUD] &
well documented \\
7 & Structural conventions (family) & not described as such & opint $=-1$
(first sample dropped at releases); plateaus as constant blocks; terminal
anchors $[0,0]$; parametrised $T_p$; 3-consonant extension; display arcs
$\cos^2(\theta/2)$; F2 offset 3 steps & polar.py: 43--51;
trajectory.py: 56; gesture.py: 487--491; config.py: 47--49 & joint sampling
jump 0.41 \Mua{} (analytic, $D=16$); display-arc gap $\approx 0.12$ [AUD] &
partial (display side in manual \S6; engine side not presented as
conventions) \\
\bottomrule
\end{tabular}
\end{adjustbox}

\medskip
\begin{minipage}{\textwidth}
\scriptsize\raggedright
Note (row 3): further registry values have \emph{no} article counterpart, so
they cannot differ from a published value --- the cluster-context angles
\code{THETA\_CC\_D} $=-2.5\pi/6$ and \code{THETA\_CC\_G} $=-\pi/12$ and the
cluster-velar radius \code{RHO\_CLUSTER\_VELAR} $=1.1$
\hbox{[R constants.py:53--55, 81]}.
\end{minipage}
\end{table}

\paragraph{Row 1/1b --- coordination matrix, Tip row (registry inheritance).}
The engine stores each row as $[\Psi_1,\Omega,\Psi_2]$
\hbox{[R constants.py:21--33]}; the Tip row is $[-2.75,\;0.5,\;\pi]$
\hbox{[R constants.py:28]}, against the article's $\Omega=0$, $\Psi_1=-3$
(Table~1) [A]: \emph{both} constants differ. Quantified by the audit: up to
0.70 \Mua{} on Tip for /d/ and /g/ (measured on the arc /d/$\to$/a/), and
zero for /b/ and /gb/, whose selection vectors exclude Tip \hbox{[AUD]}. This
row is the audit's single non-conformity, of severity below P0 \hbox{[AUD]}.
Substituting the article's Table~1 matrix into the engine (a counterfactual
run) yields identity on all seven parameters \hbox{[AUD]}, so no other row
deviates. Documentation status: omitted --- the manual asserts a different
(single-entry) difference; superseded by Corrigendum~1 (\S7).

\paragraph{Row 2 --- dental angle (registry inheritance).}
\code{THETA\_DENTAL} $=23\pi/16$ versus the article's $3\pi/2$:
$\Delta = \pi/16 = 11.25^{\circ}$
\hbox{[R constants.py:46; A \S3]}. The locus and trough demonstration scripts
explicitly restore the article value before measuring
{[R \code{scripts/}\allowbreak\code{run\_locus\_trough\_demos.py:206--216}]}, so the
published-slope demonstrations are computed at the article's angle (\S6).
This difference is absent from the manual and from code comments.

\paragraph{Row 3 --- radii (recalibration by ear; registry inheritance).}
Engine: $\rho_b = 1.2$ \hbox{[R constants.py:74]}, /i/ $=0.9$
\hbox{[R constants.py:77]}, /a/ $=0.8$ \hbox{[R constants.py:79]}; article:
$\rho_b = 1$, corner vowels $\rho = 1$ \hbox{[A \S1, \S2.1, \S3]}. The
divergence is consistent with the article's own calibration method
(``by simple hearing'' \hbox{[A \S1]}): the 2026 registry is a later state of
the same ear-tuned constants. Quantified: the manual's trough ablations
measure the /ibi/ Body peak at $+1.50$ with $\rho_b=1.2$ and $+1.25$ with the
article's $\rho_b=1.0$ (plateau $-2.25$ in both cases)
\hbox{[M l.870--872; M l.1148--1154]}. Documentation status: $\rho_b$ is
documented (the demo scripts set it explicitly and the manual discusses both
values); the vowel radii (0.9 and 0.8 vs published 1) are not flagged as
deviations --- the manual states them without comment \hbox{[M l.276--277]}.

\paragraph{Row 4 --- cluster selection (registry inheritance).}
Engine rule: for a two-consonant cluster, sum the place indices of the two
consonants (b, p, m, f, v $=0$; d, t, n, s, z $=1$; g, k $=2$; \dots) and
select $\{0,1,2,5\}$ when the sum is $\le 2$, $\{0,1,2\}$ when it is $>2$
\hbox{[R phonology.py:133--141, 337--341; constants.py:89--92]}. Executed on
the nine /b,d,g/ pairs for this supplement, the rule returns
$\{1,2,3,6\}$ (article numbering) for /bb/, /bd/, /bg/, /db/, /dd/, /gb/ and
$\{1,2,3\}$ for /dg/, /gd/, /gg/ \hbox{[R]}. The article's rule is
/b/-involvement (\S3, quote above): $\{1,2,3,6\}$ when /b/ is in the pair,
$\{1,2,3,4\}$ otherwise \hbox{[A \S3]}. Divergence: 4 of the 9 pairs
(/dd/, /dg/, /gd/, /gg/); /gb/ --- the cluster of the article's
demonstration --- conforms; among the six pairs listed in the article, /dg/
and /gd/ diverge \hbox{[AUD]}. The singleton selections do conform (/d/ and
/g/ activate $\{1,2,3,4\}$) \hbox{[R phonology.py:303--304; A \S3]}.
Documentation status: wrong in the manual --- the Demo~4 table and rule
paragraph present the article's rule as the engine's; superseded by
Corrigendum~2 (\S7).

\paragraph{Row 4a --- lineage of the cluster rule (external evidence) [EXT].}
The external note [6] compares the two original engines, \code{timit-to-Maeda}
and \code{Syllable\_Synthesis}; it never opens the article and never tests this
repository. Its findings: the isolated-consonant selectors ``are identical
after index conversion'' [EXT \S2.3]; the binary-cluster masks are identical
--- art1 rows \{0,\,1,\,2,\,5\} / \{0,\,1,\,2\} (the duplicated Jaw entry is
inert) coincide with \code{timit-to-Maeda}'s \code{ART\_CLUSTER} --- ``The
masks are therefore identical'' [EXT \S3.1]; the index rule is the same in
both engines, \code{art\_idx = 1[idx\_c1 + idx\_c2 > 2]} --- ``The decision
rule is the same'' [EXT \S3.2], with the same partition
\{bb, bd, db, bg, gb, dd\} vs.\ \{dg, gd, gg\}; the engines differ only in
trigger logic (enumerated syllable patterns vs.\ generic $m \in \{2,3\}$), in
three-consonant support, and in polar-target formulas for /d/ and /g/
[EXT \S9].

\emph{Chain of custody.} The note's \code{timit-to-Maeda} excerpts
(\code{\_CONSONANT\_ENTRIES}, \code{ART\_CLUSTER},
\code{cluster\_art\_params}) match this repository byte-for-byte at
\code{constants.py:89--92} and \code{phonology.py:337--341}, :173--180,
modulo one cosmetic character (uvular key: IPA turned-r glyph in the
repository, ASCII R in the excerpt). The engine-versus-engine identity
therefore extends to the 2026 repository at the selector layer.

\emph{Genealogical conclusion.} Because the threshold rule and both masks are
shared by the two 2023 engines, the non-article selector rule is
\emph{ancestral}: it belongs to the 2023 lineage, and the article's
/b/-conditional rule was never implemented in any engine of the family. The
2026 rework did not introduce it.

\emph{Reserve on the note.} Its two ``residual differences'' for the isolated
/d/ and /g/ polar targets are identical modulo $2\pi$
($-9\pi/16 + 2\pi = 23\pi/16$; $-\pi/12 + 2\pi = 23\pi/12$) --- a notational,
not numerical, difference [EXT \S5].

\paragraph{Row 4b --- scope of the divergence.}
\emph{Results cited in the article: none affected.}
(i)~Isolated consonants /b/, /d/, /g/ (\S3 of the article; its Figures~1
and~3): the engine's selectors conform --- /b/ $\to \{1,2,6\}$; /d/, /g/
$\to \{1,2,3,4\}$ \hbox{[R phonology.py:173--180, :303--304]}.
(ii)~The article's CC-pair text (\S3) is descriptive only: it publishes no
number and no figure for the pairs. (iii)~\S4 and Figure~4
(\code{big.bi} $\leftrightarrow$ \code{bi.gbi}): the only cluster the article
uses is /gb/, with $S_c = \{1,2,3,6\}$ --- conform, and validated end-to-end
by the audit (cluster pair /gb/: test C.4b $\le 2.2\cdot10^{-15}$; Eq.~(4)
end-to-end $\le 4.4\cdot10^{-15}$) \hbox{[AUD]}. The intersection of the
divergent pairs \{dd, dg, gd, gg\} with the article's cited results is empty.

\emph{Extended CC simulations of the supplement: affected.} Demo~5 synthesizes
six CVCCV utterances abda, abga, adba, adga, agba, agda
\hbox{[R run\_article\_demos.py:296--304]}. Exactly two run on a divergent
pair: \code{adga} (/dg/) and \code{agda} (/gd/) --- synthesized with
$S_c = \{1,2,3\}$ instead of the announced $\{1,2,3,4\}$, and their published
panel labels display the wrong vector (Corrigendum~3). All other
demonstrations (\code{big.bi}, \code{bi.gbi}, Demo~1 /ibia/, locus runs) are
untouched.

\paragraph{Row 5 --- word-onset anchoring across a pause
(author ruling 2026-10-06).}
Treated in depth in \S5. Effect: exactly 0 (bitwise) on all six published
configurations; up to 2.0812 \Mua{} outside them \hbox{[D5]} --- the largest
quantified deviation of the audit. Documentation status: partial --- README
item~4 and the in-code docstring state the rule and its motivation; the
divergence from a literal reading of the article and the conditions under
which the two readings coincide are stated for the first time in this
supplement.

\paragraph{Row 6 --- VALRECT, soft rectification (author ruling 2026-10-07).}
The article publishes no such setting; the engine applies
\begin{qbox}{\code{config.py:43--46} (comment; line breaks removed)}
\small
\code{VALRECT = 1.10 \# VLAM soft-rectification (author ruling 2026-10-07,
reduces the perceived fizz on /i/; historical canonical value from
timit-to-Maeda: 0.75; applied to the area function, not to articulatory
params)}
\end{qbox}
\noindent Quantified by the audit: locus-slope effect $\le 0.23$ (/g/ velar)
and $\le 0.014$ (/b/, /d/) \hbox{[AUD]}. Well documented: the config comment,
manual \S7.2 \hbox{[M l.721--726]} and the Demo~2 note \hbox{[M l.967--968]},
and the README headline table. One refinement: those documents call 0.75
``the article's canonical value''; per the audit the article publishes no
VALRECT --- 0.75 is the historical \code{timit-to-Maeda} default, which the
2026 ruling raises to 1.10 \hbox{[AUD; R config.py:43--46]}.

\paragraph{Row 7 --- structural conventions (family; registry inheritance).}
Not differences of published values but conventions of the Python engine that
the article does not describe as such: the first sample of each block is
dropped at releases (interpolation mode \code{opint} $=-1$)
\hbox{[R polar.py:43--51; projection.py:45--52]}; vowel and consonant
plateaus are materialized as constant blocks (\code{np.tile})
\hbox{[R trajectory.py:56]}; utterance edges carry terminal anchors $[0,0]$
\hbox{[R gesture.py:487--491]}; pause durations $T_p$ are parametrised; a
three-consonant cluster extension exists with no article counterpart
\hbox{[R phonology.py:342--352]}; the polar \emph{display} draws arcs with a
$\cos^2(\theta/2)$ profile while the engine uses $\cos(\theta/2)$
\hbox{[R polar\_primitives.py; M \S6]}; and the locus measurement samples F2
30~ms after release (\code{LOCUS\_OFFSET\_STEPS} $=3$)
\hbox{[R config.py:47--49]}. Joint effect of the sampling-side conventions:
an analytic 0.41 \Mua{} jump at block joints ($D=16$); display-vs-engine arc
gap $\approx 0.12$ \Mua{} \hbox{[AUD]}. Documentation status: partial --- the
display conventions are documented (manual \S6, ``Polar Trajectory
Reconstruction''); the engine-side conventions (plateau blocks, \code{opint},
the 0.41 \Mua{} jump) are not presented as conventions.

\medskip
\noindent\textbf{Ranking at trajectory level.} Outside the published grid,
the differences rank: word-onset rule 2.08 $>$ CO Tip row 0.70 $>$ joint
sampling jump 0.41 $>$ display-arc gap 0.12 \Mua{} \hbox{[AUD; D5]}. Inside
the published grid, rows 1, 3, 6 and 7 act at their quantified magnitudes
while row~5 is exactly zero (\S5).

\section{Focus: word-onset anchoring across a pause}

\textbf{The two readings.} The article defines the pause as a diphthong
transition between the previous $V_e$ and the next $V_o$ (Q2), with
$V_o = (\delta_o\,\rho_V, \theta_V)$ anchored on ``the closest vowel'' (Q1)
\hbox{[A \S2.2]}. For a non-initial word that begins with a consonant after a
pause, two readings of ``closest vowel'' --- and of the scaling coefficient
--- are possible:
\begin{center}
\begin{minipage}[t]{0.485\textwidth}
\begin{qbox}{Engine rule (2026) [R gesture.py:385--432]}
\small
$V_{o,2} \;:=\; (\delta_e \cdot \rho_{V_1},\; \theta_{V_1})$
\\[3pt]
The onset anchors on the \emph{previous} word's $V_e$ point, scaled by the
\emph{end} coefficient $\delta_e$.
\\[3pt]
Author-validated 2026-10-06 (docstring Q6 below).
\end{qbox}
\end{minipage}\hfill
\begin{minipage}[t]{0.485\textwidth}
\begin{qbox}{Article-literal reading (2023) [A \S2.2]}
\small
$V_{o,2} \;=\; (\delta_o \cdot \rho_{V_2},\; \theta_{V_2})$
\\[3pt]
The onset anchors on the word's \emph{own} closest vowel, scaled by the
\emph{onset} coefficient $\delta_o$, as at an utterance-initial word.
\end{qbox}
\end{minipage}
\end{center}

\begin{qbox}{\code{synthSYL/gesture.py:399--406} --- docstring (Q6, verbatim)}
\small
``Non-initial word (after a pause): the pause is the diphthong transition
between the PREVIOUS $V_e$ and the next $V_o$ (article \S2.2), so the onset
anchors on the previous word's $V_e$ point --- $V_o := (\delta_e \cdot
\rho_{\text{prev}}, \theta_{\text{prev}})$. Consequences: with
$\text{COEFCEN} = 1$ ($\delta_e = 1$, $V_e = V$) the word-2 onset departs
from the FULL vowel and there is no return to the half-radius schwa at the
word boundary; with the symmetric $\delta_o = \delta_e$ the value is
unchanged (0.45 for /i/), preserving the `big@bi' schwa demonstrations.
(Author-validated 2026-10-06.)''
\end{qbox}

\textbf{Coincidence lemma.} The two readings produce the same anchor if and
only if
\begin{equation}
\delta_e\,\rho_{V_1} \;=\; \delta_o\,\rho_{V_2}
\qquad\text{and}\qquad
\theta_{V_1} \equiv \theta_{V_2} \pmod{2\pi}.
\label{eq:coincidence}
\end{equation}
\emph{Corollary.} At the published conditions --- the article always sets
$\delta_o=\delta_e=\delta$, and the demonstrations place the same vowel on
both sides of the pause --- both conditions hold identically, so the two
readings coincide exactly. For \code{bi\,bi} with the engine's /i/
($\rho=0.9$, $\theta=5\pi/3$): both readings give $V_{o,2} = (0.45,
5\pi/3)$ at $\delta=0.5$, $(0.63, 5\pi/3)$ at $\delta=0.7$, $(0.90,
5\pi/3)$ at $\delta=1.0$ \hbox{[R; M \S4]}.

\textbf{Structural argument.} The single-word utterances \code{big.bi} and
\code{bi.gbi} contain no non-initial word onset at all: their internal `.`
boundary is an inter-syllable boundary governed by the chaining rule (README
item~2, the \code{COEFCEN}-weighted `.' C.C anchor), an assembly rule
identical under both readings. The number of divergent application sites in
these demonstrations is therefore zero for structural reasons, independently
of the lemma \hbox{[D5]}.

\textbf{The D5 experiment.} Each case was run twice through the complete
pipeline, identical except for \code{\_insert\_vowel\_onset\_anchors}
(engine rule vs literal rule); the table reports the maximum absolute
difference of the 7-parameter flow \hbox{[D5]}:

\begin{table}[htbp]
\centering
\caption{Difference-5 experiment: engine rule vs article-literal rule at the
published conditions ($\delta_o=\delta_e=\delta$) and at one control outside
them. max $|\Delta P|$ over the 7 Maeda parameters and the whole utterance
[D5].}
\label{tab:d5}
\footnotesize
\begin{adjustbox}{max width=\textwidth,center}
\begin{tabular}{@{}>{\raggedright\arraybackslash}p{2.7cm}
                >{\raggedright\arraybackslash}p{1.9cm}
                >{\raggedright\arraybackslash}p{1.55cm}
                >{\raggedright\arraybackslash}p{2.1cm}
                >{\raggedright\arraybackslash}p{2.6cm}
                >{\raggedright\arraybackslash}p{4.1cm}@{}}
\toprule
\textbf{Utterance} & $\boldsymbol\delta$ & \textbf{Onset sites} &
\textbf{Divergent branch} & \textbf{max $|\Delta P|$ (\Mua)} &
\textbf{Reading} \\
\midrule
bi bi & 0.5 & 2 & 1 (non-initial word) & 0 (exact, bitwise) &
coincidence (lemma~\eqref{eq:coincidence}) \\
bi bi & 0.7 & 2 & 1 (non-initial word) & 0 (exact) & coincidence (lemma) \\
bi bi & 1.0 & 2 & 1 (non-initial word) & 0 (exact) & coincidence (lemma) \\
big.bi & 0.7 & 1 & 0 & 0 (exact) & structural (no site) \\
big.bi & 1.0 & 1 & 0 & 0 (exact) & structural (no site) \\
bi.gbi & 1.0 & 1 & 0 & 0 (exact) & structural (no site) \\
control: ``big ba'' & $\delta_o{=}0.5$, $\delta_e{=}0.7$ & 2 & 1 &
\textbf{2.0812} & divergence \\
\bottomrule
\end{tabular}
\end{adjustbox}
\end{table}

\textbf{The control in detail.} For \code{big\,ba} with $\delta_o = 0.5$,
$\delta_e = 0.7$, the two anchors split: the engine rule gives $(0.63,
5\pi/3)$ --- $\delta_e \cdot 0.9$ on the previous /i/ --- while the literal
rule gives $(0.40, \pi)$ --- $\delta_o \cdot 0.8$ on the word's own /a/; the
polar gap is $\Delta\rho = 0.23$ and 1.73~rad on the unit circle (matching
audit test D.6 \hbox{[AUD]}). Per-parameter maxima: Body 2.0812 (at
$t = 129$ steps $\approx 1.29$~s), Tip 1.99, LipH 1.81, Jaw 1.08, while
Dorsum, LipP and Hy remain at 0.23--0.35; word~1 is untouched
($\Delta P = 0$ for $t < 112$ steps) \hbox{[D5]}. The four strongly affected
parameters are the consonant selection of the /b/, $S_c = \{1,2,6\}$ = Jaw,
Body, LipH, together with Tip from the vowel selection $S_v = \{3,4,5,7\}$ =
Dorsum, Tip, LipP, Hy --- consistent with the split of article Eq.~4
\hbox{[D5]}.

\textbf{Interpretation.} Within the published grid the rule is invisible:
Table~\ref{tab:d5} shows exact zeros for all six published configurations,
and the README records that ``all reference reproductions remain
bit-identical'' (Q10, quoted in \S6). The rule bites only outside that grid
--- an asymmetric $\delta_o \ne \delta_e$, or a different vowel on the two
sides of a pause. Users who explore such conditions should be aware that the
engine implements the previous-Ve anchoring, not the literal reading.

\section{Why the published demonstrations are nevertheless reproduced}

The audit's evidence, case by case \hbox{[AUD]}:
\begin{itemize}\itemsep2pt\parsep0pt
\item \textbf{Equation level.} Independent implementations of Eqs.~(1)--(4)
agree with the engine to $\le 5\cdot10^{-15}$.
\item \textbf{Counterfactual CO.} Substituting the article's Table~1
coordination matrix into the engine yields identity to $\le 10^{-15}$ on all
seven parameters: the Tip-row deviation (row~1/1b) is a pure parameter
substitution, not an equation change.
\item \textbf{End-to-end transformation.} The \code{big.bi} $\to$
\code{bi.gbi} transformation is reconstructed from the article's equations
alone to $\le 4.4\cdot10^{-15}$ on six parameters, with the published gain of
one period $T = 160$~ms (176 $\to$ 160 steps).
\item \textbf{Locus demonstration.} With the article's positions restored
--- which the repository's own scripts do
{[R \code{scripts/}\allowbreak\code{run\_locus\_trough\_demos.py:206--216}]} --- the
measured locus slopes are 0.7028 (/b/), 0.4211 (/d/), 2.3352 (/g/ velar) and
0.6305 (/g/ palatal), reproduced to $\le 5\cdot10^{-4}$. These are the values
reported in the README headline table and in manual \S10.2
\hbox{[M l.959--962]}.
\item \textbf{Word-onset rule.} Exactly zero difference at all published
conditions \hbox{[D5]} (\S5, Table~\ref{tab:d5}).
\end{itemize}

\begin{qbox}{README, ``Key Modifications'', item 4 --- closing sentence
(Q10, verbatim [R])}
\small
``\dots with the symmetric $\delta_o = \delta_e$ the value is unchanged (the
`big@bi' schwa demonstrations are preserved, and all reference reproductions
remain bit-identical).''
\end{qbox}

\noindent The audit tally: 104 tests --- 91 conforming, 10 documented
deviations (the families of \S4), 1 non-conformity (the CO Tip row), and 0
P0-severity issues \hbox{[AUD]}. \textbf{Consequence:} the engine faithfully
applies the article's equations; every deviation is a parameter or rule
choice. No engine modification is warranted; the appropriate fix is
documentation --- this supplement and the two corrigenda of \S7.

\section{Corrigenda to the manual}

The audit found exactly two factually wrong statements in the manual
\hbox{[AUD]}, plus one family of wrong labels in the demonstration artifacts
generated by the repository's own figure generator; the three corrections are
issued below as Corrigenda~1--3. Each corrigendum gives (a) the location, (b)
the current text or state quoted verbatim, (c) a ready-to-paste corrected
text, and (d) the verified engine values. \emph{Application status (October
14, 2026):} the three corrections have been applied to the manual and the
figure labels regenerated from the engine; the quoted ``current'' texts
refer to the state at commit \code{4a0d76290aee}.

\subsection*{Corrigendum 1 --- coordination-matrix claim
(manual \S1.2, lines 225--227)}

\textbf{(b) Current text [M l.225--227], verbatim:}
\begin{qbox}{\code{docs/manual.tex} lines 225--227 (Q8, verbatim)}
\small
``The CO matrix matches the article's Table~1 almost exactly; the only
difference is the Tip parameter $c_0$: $-2.75$ in
\code{synthSYL/constants.py} vs $-3.0$ in the article.''
\end{qbox}
\noindent This statement is wrong on two counts: the Tip row differs in
\emph{two} constants, not one; and no other row differs.

\textbf{(c) Ready-to-paste corrected text:}
\begin{pastebox}{Suggested replacement for manual \S1.2, lines 225--227}
\small
``The CO matrix matches the article's Table~1 except for the Tip row, where
two parameters differ: $\Omega = 0.5$ (article: 0) and $\Psi_1 = -2.75$
(article: $-3.0$) in \code{synthSYL/constants.py}; all other rows are
identical. The deviation is inherited from the original timit-to-Maeda
registry and is quantified in the Differential Supplement \S4 (up to 0.70
Maeda units on Tip for /d/, /g/; zero for /b/ and /gb/).''
\end{pastebox}

\noindent\textbf{(d) Verified engine values.}
\code{constants.py:28} stores the Tip row as $[\Psi_1, \Omega, \Psi_2] =
[-2.75,\;0.5,\;\pi]$ \hbox{[R]}; the article's Table~1 gives $\Omega = 0$ and
$\Psi_1 = -3$ \hbox{[A]}; effect up to 0.70 \Mua{} on Tip for /d/, /g/, zero
for /b/ and /gb/ \hbox{[AUD]}.

\subsection*{Corrigendum 2 --- cluster-selection rule
(manual \S10.4, table lines 1260--1273 and rule paragraph lines 1275--1280)}

\textbf{(b) Current text [M l.1275--1280], verbatim:}
\begin{qbox}{\code{docs/manual.tex} lines 1275--1280 (Q9, verbatim)}
\small
``Rule: when /b/ is involved in the pair, the selection vector is
$S_c = \{1,2,3,6\}$ (article numbering; code $\{0,1,2,5\}$ = jaw, body,
dorsum, LipH); otherwise $S_c = \{1,2,3,4\}$ (article numbering; code
$\{0,1,2,3\}$ = jaw, body, dorsum, tip). The /g/ is systematically palatal
when in a cluster.''
\end{qbox}
\noindent This is the \emph{article's} rule (\S3, quote in \S3 above),
presented as the engine's. The Demo~4 table cells are consistent with it and
therefore also wrong for /dg/ and /gd/ ($\{1,2,3,4\}$ is a vector the engine
never returns for a cluster; the engine returns $\{1,2,3\}$ for both).
\textbf{Explicit cell replacements}: /dg/ row
\hbox{[M l.1268]} --- $S_c$ \{1,\,2,\,3,\,4\} $\to$ \{1,\,2,\,3\} and note
``no /b/ $\to S_c{=}\{1,2,3,4\}$'' $\to$ ``engine threshold rule
(place-index sum $> 2$)''; /gd/ row \hbox{[M l.1270]} likewise.

\textbf{(c) Ready-to-paste corrected text:}
\begin{pastebox}{Suggested replacement for the manual \S10.4 rule paragraph}
\small
``Engine rule (verified): the selection vector is $S_c = \{1,2,3,6\}$ (code
$\{0,1,2,5\}$) when the sum of place indices ($b = 0$, $d = 1$, $g = 2$) is
$\le 2$, and $S_c = \{1,2,3\}$ (code $\{0,1,2\}$) otherwise. Hence /bb/,
/bd/, /bg/, /db/, /dd/, /gb/ use $\{1,2,3,6\}$ and /dg/, /gd/, /gg/ use
$\{1,2,3\}$; the code vector $\{0,1,2,3\}$ is never produced for a cluster
(it is the singleton selection of /d/ and /g/). The article's rule ---
$\{1,2,3,6\}$ for clusters containing /b/ (e.g.\ /gb/), $\{1,2,3,4\}$
otherwise --- differs from the engine's for the four pairs /dd/, /dg/, /gd/,
/gg/. The engine's rule is inherited from the 2023 lineage: both original
engines implement the same threshold rule [EXT]; the article's
/b/-conditional rule was never implemented in the family. See Differential
Supplement \S4, rows~4 and~4a.''
\end{pastebox}

\noindent\textbf{(d) Verified engine values.} Place indices $b{=}0$,
$d{=}1$, $g{=}2$ \hbox{[R phonology.py:133--141]}; the rule selects
\code{ART\_CLUSTER[0]} $=\{0,1,2,5\}$ when the sum of the two indices is
$\le 2$ and \code{ART\_CLUSTER[1]} $=\{0,1,2\}$ otherwise
\hbox{[R phonology.py:337--341; constants.py:89--92]}, executed on the nine
pairs for this supplement: \{0,1,2,5\} for bb, bd, bg, db, dd, gb;
\{0,1,2\} for dg, gd, gg (article numbering: $\{1,2,3,6\}$ and
$\{1,2,3\}$) \hbox{[R]}. The article's rule is quoted verbatim in \S3
\hbox{[A \S3]}. The affected utterances of the extended simulations are
\code{adga} and \code{agda} (see Row~4b in \S4).

\subsection*{Corrigendum 3 --- demonstration-figure labels
(generator \code{run\_article\_demos.py}; manual embeds at lines 1284
and 1291)}

\textbf{(a) Location.} The demonstration figure generator
\code{scripts/run\_article\_demos.py} hardcodes the article's labels in
\code{CLUSTER\_INFO} (lines 91--98): the /dg/ and /gd/ entries carry
\code{"Sc": "\{1,2,3,4\}"} with note \code{"no /b/ -> Sc=\{1,2,3,4\}"} (verbatim
excerpt in the brief's Annex~A). These strings flow into (i)~the console log
(line 310), (ii)~the panel titles of \code{cluster\_params.png} (line 356),
and (iii)~the panel titles of \code{cluster\_polar.png} (line 416). The manual
embeds both PNGs (\hbox{[M l.1284]} and \hbox{[M l.1291]}), so the printed
Demo~4 section inherits the mislabels.

\textbf{(b) Current state.} The synthesis that produced those very panels ran
with $S_c = \{1,2,3\}$ for /dg/ and /gd/ (threshold rule; \S4, row~4, verified
by execution). The published artifacts therefore display $S_c = \{1,2,3,4\}$
--- a vector the engine never returns for a cluster --- next to trajectories
that are faithful to the engine: a label/engine inconsistency inside the
published Demo~4 artifacts. Only the labels are wrong. The spectrogram titles
carry no $S_c$ label, and Demo~1 (/ibia/) is unaffected.

\textbf{(c) Correction --- two coordinated parts.}
\emph{Generator (maintainer action).} Option~B (recommended): derive the
labels at runtime from \code{synthSYL.phonology.cluster\_art\_params},
converting code indices to article numbering by $+1$, so the displayed labels
can never drift from the engine again; keep a static note field for prose
only. Option~A (minimal): edit the two entries to \code{"Sc": "\{1,2,3\}"} with
note ``engine threshold rule (place-index sum $> 2$)''. Then re-run the
generator, propagate the regenerated \code{cluster\_params.png} and
\code{cluster\_polar.png} into \code{docs/figures/demo5/}, and rebuild the
manual PDF. \emph{Ready-to-paste manual note (immediate documentation fix),
to insert under the Demo~4 figures:}
\begin{pastebox}{Suggested note under the manual Demo~4 figures}
\small
``Labels note. In the figures above, the $S_c$ annotations for /dg/ and /gd/
($S_c = \{1,2,3,4\}$) are legacy labels: the synthesis runs with
$S_c = \{1,2,3\}$ (threshold rule on the place-index sum; see the Differential
Supplement, \S4 row~4, and Corrigendum~2). The trajectories are unaffected;
corrected figures are regenerated by \code{scripts/run\_article\_demos.py}.''
\end{pastebox}

\textbf{(d) Verified engine values.} /dg/ and /gd/ $\to$ code $[0,1,2]
\equiv \{1,2,3\}$ (article numbering); the four panels /bd/, /bg/, /db/,
/gb/ already display the correct $\{1,2,3,6\}$. Panel by panel:

\begin{center}
\footnotesize
\begin{tabular}{@{}lccc@{}}
\toprule
\textbf{Panel} & \textbf{Displayed label} & \textbf{Engine value} &
\textbf{Action} \\
\midrule
/bd/, /bg/, /db/, /gb/ & $S_c = \{1,2,3,6\}$ & $\{1,2,3,6\}$ & none \\
/dg/ & $S_c = \{1,2,3,4\}$ + ``no /b/'' note & $\{1,2,3\}$ & relabel \\
/gd/ & $S_c = \{1,2,3,4\}$ + ``no /b/'' note & $\{1,2,3\}$ & relabel \\
\bottomrule
\end{tabular}
\end{center}

\section{Version and provenance statement}

The following block is ready to paste into the README and the manual preface:
\begin{pastebox}{Version and provenance statement (ready to paste)}
\small\raggedright
\textbf{Version and provenance.} The demonstrations in this repository were
produced with the Python engine synthSYL reworked in 2026 within
big.bi-supplement (commit \code{4a0d76290aee},
DOI 10.5281/zenodo.23239483), which extends the author's timit-to-Maeda
package. The 2023 article (arXiv:2307.02299) reports demonstrations produced
with the author's MATLAB implementation of the same model. The registry
parameters are inherited from the engine and recalibrated by ear; they do not
all coincide with the values published in the 2023 article --- the complete
differential is catalogued in the Implementation Differential and Provenance
Supplement (\code{docs/differential\_supplement.pdf}, \S4), and none of the
differences is an equation-level change. Two choices are dated author
rulings: the word-onset anchoring rule (2026-10-06) and
$\mathrm{VALRECT} = 1.10$ (2026-10-07).
\end{pastebox}

\section{Reproducibility}

\textbf{Environment.} Python $\ge$ 3.10 (the audit ran 3.12) with
numpy, scipy, matplotlib and pytest; repository at commit
\code{4a0d76290aee}; MIT license \hbox{[R]}. The commands below follow the
audit workspace \hbox{[AUD; D5]}; steps 3--5 use the scripts
\code{exp\_d5\_onset\_rule.py} and \code{fig\_d5.py} provided there (on
Windows, replace \code{venv/bin/} with \code{venv\textbackslash Scripts\textbackslash} and set
\code{BIGBI\_REPO} with \code{set}):
\begin{quote}\small
\begin{verbatim}
# 1. environment
python -m venv venv && venv/bin/pip install numpy scipy matplotlib pytest
# 2. repository at the pinned commit
git clone https://github.com/FBerthommier/big.bi-supplement
cd big.bi-supplement && git checkout 4a0d76290aee
# 3. difference-5 experiment (~10 s)   [script: audit workspace]
BIGBI_REPO=/path/to/repo python exp_d5_onset_rule.py
#    -> d5_results.json + d5_<case>_{code,lit}.csv (14 CSVs)
# 4. figures
python fig_d5.py   # -> fig_d5_bibi.png, fig_d5_bigbi.png, fig_d5_synthesis.png
# 5. full 104-test audit suite (~10 min)
cd tests/compliance && BIGBI_REPO=/path/to/repo python -m pytest . -q
\end{verbatim}
\end{quote}

\textbf{Artifact inventory.} Difference-5: \code{d5\_results.json}, 14
per-case CSVs (engine vs literal rule), three figures
(\code{fig\_d5\_bibi.png}, \code{fig\_d5\_bigbi.png},
\code{fig\_d5\_synthesis.png}) \hbox{[D5]}. Audit: the 104-test matrix
(\code{compliance\_report.md}), the engine-vs-article correspondence table
(\code{annexe\_correspondance\_}\allowbreak\code{synthsyl\_article.md}), and
the 13-page difference-5 report
(\code{difference5\_effect\_audit.pdf}) \hbox{[AUD]}.

\textbf{Selector probe.} Every selector claim of this document is
re-derivable from the repository alone. Run from any directory with the
repository root in \code{sys.path}:
\begin{quote}\small
\begin{verbatim}
import sys
sys.path.insert(0, "<repository root>")
from synthSYL.phonology import cluster_art_params

def article_rule(pair):          # article S3, p. 4
    return {1, 2, 3, 6} if "b" in pair else {1, 2, 3, 4}

for p in ["bb", "bd", "bg", "db", "dd", "dg", "gb", "gd", "gg"]:
    code = cluster_art_params(list(p))          # 0-based code indices
    art  = sorted(i + 1 for i in code)          # article numbering (+1)
    ref  = sorted(article_rule(p))
    print(f"/{p}/  engine {art}  article {ref}  "
          f"{'conform' if art == ref else 'DIVERGES'}")
\end{verbatim}
\end{quote}
\noindent Expected output: \code{conform} for bb, bd, bg, db, gb (engine
$\{1,2,3,6\}$); \code{DIVERGES} for dd (engine $\{1,2,3,6\}$) and for dg, gd,
gg (engine $\{1,2,3\}$); the article rule answers $\{1,2,3,4\}$ for these
four. Place-index sums ($b{=}0$, $d{=}1$, $g{=}2$): bb 0; bd 1; bg 2; db 1;
dd 2; dg 3; gb 2; gd 3; gg 4 --- the threshold is sum $> 2$.

\begin{thebibliography}{9}\small

\bibitem{article}
F.~Berthommier, ``Why can big.bi be changed to bi.gbi? A mathematical model
of syllabification and articulatory synthesis'', arXiv:2307.02299 [eess.AS],
July 2023 (Interspeech 2023). Cited as [A].

\bibitem{repo}
F.~Berthommier, \emph{big.bi-supplement} repository, commit
\code{4a0d76290aee}, MIT license, DOI 10.5281/zenodo.23239483, 2026.
\url{https://github.com/FBerthommier/big.bi-supplement}. Cited as [R].

\bibitem{manual}
F.~Berthommier, \emph{Didactic Manual: big.bi $\to$ bi.gbi articulatory
simulation}, \code{docs/manual.pdf}, 26 pages, at the same commit. Cited as
[M].

\bibitem{audit}
104-test compliance audit report, \code{compliance\_report.md}, October 2026.
Cited as [AUD].

\bibitem{d5}
Difference-5 effect report, \code{difference5\_effect\_audit.pdf}, 13 pages,
October 2026. Cited as [D5].

\bibitem{ext}
Anonymous, ``Comparative audit of consonant selection vectors
$C \in \{$b, d, g$\}$ in Timit-to-Maeda and Syllable\_Synthesis,''
technical synthesis for verification of reasoning, 6 pp., October 13, 2026
(file \code{selectors\_CC\_comparison.pdf}). Cited as [EXT].

\end{thebibliography}

\end{document}
